\documentclass[11pt,letterpaper]{article}
\usepackage[T1]{fontenc}
\usepackage[utf8]{inputenc}
\usepackage{lmodern,microtype}
\usepackage[margin=1in]{geometry}
\usepackage{amsmath,amssymb,amsthm,mathtools,stmaryrd}
\usepackage{booktabs,longtable,array}
\usepackage[dvipsnames]{xcolor}
\usepackage{listings}
\usepackage[most]{tcolorbox}
\usepackage{enumitem}
\usepackage{fancyhdr}
\usepackage{hyperref}
\hypersetup{colorlinks=true,linkcolor=MidnightBlue,citecolor=MidnightBlue,urlcolor=MidnightBlue,
 pdftitle={Beyond BusyLean's Contender 6: Autonomous Model Ramification},
 pdfauthor={Research proposal prepared for Vladimir Reshetnikov},
 pdfsubject={Computable large numbers, model ramification, Lean 4, exact domination}}
\setlength{\headheight}{14pt}
\pagestyle{fancy}
\fancyhf{}
\fancyhead[L]{\small Beyond BusyLean's Contender 6}
\fancyhead[R]{\small Research proposal}
\fancyfoot[C]{\thepage}
\setlength{\parskip}{4pt}
\setlength{\parindent}{1.3em}
\setcounter{tocdepth}{1}
\emergencystretch=2em
\newtheorem{theorem}{Theorem}[section]
\newtheorem{lemma}[theorem]{Lemma}
\newtheorem{proposition}[theorem]{Proposition}
\newtheorem{corollary}[theorem]{Corollary}
\theoremstyle{definition}
\newtheorem{definition}[theorem]{Definition}
\newtheorem{remark}[theorem]{Remark}
\newcommand{\N}{\mathbb N}
\newcommand{\W}{\mathcal W}
\newcommand{\Fam}{\mathcal F}
\newcommand{\Mod}{\mathcal M}
\newcommand{\Tree}{\mathsf{Tree}}
\newcommand{\Eval}{\operatorname{eval}}
\newcommand{\rank}{\operatorname{rank}}
\newcommand{\Next}{\operatorname{Next}}
\newcommand{\MJoin}{\operatorname{Join}}
\newcommand{\Ram}{\operatorname{Ram}}
\newcommand{\Val}{\operatorname{C}}
\newcommand{\idx}{\operatorname{I}}
\newcommand{\td}{\operatorname{treeOf}}
\newcommand{\ps}{\mathsf p}
\newcommand{\FF}{\operatorname{F}}
\newcommand{\Iter}[2]{#1^{[#2]}}
\newcommand{\code}[1]{\texttt{#1}}
\newcommand{\src}[1]{\texttt{\detokenize{#1}}}
\newtcolorbox{statusbox}{colback=gray!5,colframe=gray!55,boxrule=0.5pt,
 left=9pt,right=9pt,top=7pt,bottom=7pt,breakable}
\lstdefinestyle{lean}{basicstyle=\ttfamily\footnotesize,columns=fullflexible,
 keepspaces=true,breaklines=true,showstringspaces=false,
 commentstyle=\color{gray},keywordstyle=\color{MidnightBlue},
 morekeywords={def,theorem,structure,where,namespace,end,import,abbrev,fun,match,with,
 let,in,by,exact,have,inductive,termination_by,structural},
 morecomment=[l]{--},morecomment=[s]{/-}{-/},
 literate={→}{{$\to$}}1 {←}{{$\leftarrow$}}1 {≤}{{$\le$}}1 {≥}{{$\ge$}}1
 {∀}{{$\forall$}}1 {∃}{{$\exists$}}1 {⟨}{{$\langle$}}1 {⟩}{{$\rangle$}}1
 {ℕ}{{$\N$}}1 {¬}{{$\neg$}}1 {∈}{{$\in$}}1 {×}{{$\times$}}1 {α}{{$\alpha$}}1}
\lstset{style=lean}

\begin{document}
\begin{center}
{\Large\bfseries Beyond BusyLean's Contender 6}\\[6pt]
{\large Autonomous Model Ramification and a\par
46-fold Tower-Feedback Separation}\\[12pt]
Research proposal prepared for Vladimir Reshetnikov\\
October 12, 2026
\end{center}

\begin{abstract}
We propose a closed, computable Lean definition that strengthens the current
BusyLean contender by changing what the object language can interpret, not
merely by enlarging its numerical argument. A fixed seven-slot interface makes
lower reflection operators, forest recursors, and program-generated tree
schedules available to higher interpretations. The interface closes under all
finite model ramifications, and is then ramified structurally along a tree
produced by the strengthened language itself. A depth-11 program provides a
uniform comparison certificate, valid for every external oracle. If $C_6$ is
the adopted contender at the pinned repository revision, and
$\Theta(x)=(y\mapsto 2^y)^{[x+1]}(x+1)$, the construction satisfies
\[
                    \boxed{\Iter{\Theta}{46}(C_6)<C_7.}
\]
We also prove an interpretation-relative cofinality theorem for model
ramification along definable trees, and a diagonal non-majorization theorem.
These results require no assigned proof-theoretic ordinal. The package
contains detailed mathematical proofs, new Lean definitions and 84 theorem
scripts, exact baseline source copies, and a guarded reference interpreter.
The new Lean sources have not been compiler-checked in this environment;
1,830 finite reference-model assertions passed, but kernel certification,
axiom closure, and contest build-time compliance remain unverified.
\end{abstract}

\begin{statusbox}
\textbf{Validation boundary.} This is a mathematical research proposal with an
unverified Lean formalization, not an already certified contest submission.
No Lean compiler was available during preparation, and retrieval attempts did
not supply one. The recorded Lean status is \code{not\_run}. There are no
intentional proof placeholders in the new sources, but elaboration or proof
repairs may still be necessary. Neither contender has been evaluated.
\end{statusbox}

\noindent\textbf{Pinned baseline.} All comparisons refer to BusyLean commit\\
\hspace*{1em}\code{86b360648163f0249eab79539438b6a83dfe8215}.\\
\textbf{Target toolchain.} Lean 4.34.0; bundled Lean libraries only; no Mathlib.

\newpage
\tableofcontents
\newpage

\section{The result and its intended scope}
\label{sec:result}

Let $\Iter f0(x)=x$ and $\Iter f{k+1}(x)=f(\Iter fk(x))$. Throughout this article,
square brackets on an exponent denote \emph{function iteration}; they never
denote ordinary exponentiation. Define
\begin{equation}\label{eq:theta}
 E(x)=2^x,\qquad \Theta(x)=\Iter E{x+1}(x+1).
\end{equation}
Thus $\Theta(0)=2$, $\Theta(1)=16$, and $\Theta(2)=2^{256}$.
For general $x$, the input determines both the initial exponent and the number
of exponentiation steps. In particular, $\Theta(C_6)$ is already much more
than an exponential or a fixed-height tower in $C_6$.

\begin{theorem}[Main numerical separation]\label{thm:main}
The construction in Section~\ref{sec:models} defines a natural number $C_7$
using total interpretation and structural recursion. With $C_6$ the adopted
BusyLean number at the pinned revision,
\[
                  \Iter\Theta{46}(C_6)<C_7.
\]
Consequently $C_6<C_7$.
\end{theorem}

The proof is the specialization at $n=42$ of the stronger uniform statement
\begin{equation}\label{eq:uniform-main}
 \forall n\ge 11,\qquad
 \Iter\Theta{n+4}\bigl(\FF_6(n)\bigr)<\Val(M_{\mathrm{final}},n),
\end{equation}
where $\FF_6$ is the repository's \code{BusyLean.finalFamily}. The value
$46=42+4$ is an explicitly proved conservative lower bound. It is not an
estimate of the entire growth strength of the construction.

The principal advance is structural. The repository's adopted construction
runs a fixed reflection mechanism along a program-generated tree. The new
construction makes the mechanism itself available to subsequent interpreters,
retains complete families of lower mechanisms, and performs
\emph{model-valued} recursion along a generated tree. This changes the
operations that a short program can call. The numerical certificate in
Theorem~\ref{thm:main} is a transparent way to verify a large gap without
assigning ordinals to these operations.

There is no claim of an absolutely largest computable number, of an optimal
use of the contest time budget, or of a newly established proof-theoretic
ordinal. The direct scalar comparison is with the adopted $C_6$. The
repository's established $C_5<C_6$ and proposal-P5 comparison yield further
corollaries once the new comparison is kernel-checked and linked against the
full repository \cite{repo-comparison}. Exact comparisons with the separate P1--P4 proposal
constants are \emph{not} established here. Their heuristic ranking is not a
substitute for such proofs.

\section{What is being extended}
\label{sec:baseline}

\subsection{The current repository, not an earlier proposed contender}

The pinned root README and source identify the adopted number as
\begin{equation}\label{eq:old}
 C_6=\FF_6(42),\qquad
 \FF_6(n)=S_{\mathrm{old}}(\ps,I_{\mathrm{old}},n).
\end{equation}
Here $S_{\mathrm{old}}$ is forest reflection over the finite-type W-language,
$\ps$ is P5's proposed number as a family in its input, and
$I_{\mathrm{old}}$ joins the W-trees denoted by closed programs of that
language. These are the actual adopted definitions, not the earlier P5
constant also called \code{contender\_6} in its own namespace
\cite{repo-readme,repo-c6,repo-reflection,p5-defs}.

The inherited seed has type $\Fam=\N\to\N\to\N$ and is
\[
 \ps(i,n)=\code{BigNumber.iterate}\;
           \code{BigNumber.diagonal}\;n\;n.
\]
It ignores its first argument. The full definition is retained byte-for-byte
in the vendored baseline; it includes P5's seed guard over the old numerical
maximum. We do not claim to have eliminated that inherited dependence.

The old reflector is
\begin{equation}\label{eq:old-reflection}
 R_{\mathrm{old}}(f,n)=n+1+
 \max\left\{\max_{i\le n}f(i,n),\
      \max_{t\in\mathcal T_n}\Eval_{\mathrm{old},f}(t)\right\},
\end{equation}
where $\mathcal T_n$ is the repository's finite raw-term enumeration. We use
$\max\varnothing=0$ for numerical maxima throughout. At a
leaf, the old forest uses this reflector with the seed; at a node, it uses the
family of child stages. Our new forest has the same outer recursive shape,
but a strengthened interpretation in its finite maximum.

\subsection{Contest obligations and what remains outstanding}

The repository rules demand a closed computable natural, a formal strict
comparison, no prohibited computational shortcuts or axioms, and clean-build
budgets of 15 seconds for definitions and 60 seconds for proofs. They prohibit
using the previous named contender as a shortcut and allow only
\code{propext} and \code{Quot.sound} among axioms. The pinned toolchain is
Lean 4.34.0 \cite{repo-readme,lean-release}.

The proposed definition never names \code{BusyLean.contender\_6},
\code{BusyLean.finalFamily}, or \code{BigNumber.contender\_6}. It deliberately
reuses their generic constituents: the older reflector, forest operator,
index, and P5 seed family. A short program can reconstruct the old final
family from those constituents, and this is exactly the comparison witness.
Thus the proposal offers syntactic non-use of the old scalar together with a
substantial general extension; it does \emph{not} pretend to be semantically
independent of the old construction.

The mathematical comparison is proved below. The supplied Lean proof scripts
are a concrete formalization attempt, not a certified result. Until the clean
build, axiom audit, and dependency audit run successfully, the formal-proof
and timing obligations remain open engineering checks for this submission.
No elapsed build times are inferred from the size of the sources.

\subsection{Reproducible baseline}

Three baseline files suffice for the standalone proposed definitions. They
were reconstructed from the pinned GitHub contents and checked against their
Git blob identifiers:
\begin{center}
\begin{tabular}{@{}lr@{}}
\toprule
Vendored path & Bytes\\
\midrule
\src{lean/BigNumber/Definitions.lean} & 6,834\\
\src{lean/BusyLean/Reflection.lean} & 2,841\\
\src{lean/BusyLean/Contender6.lean} & 2,033\\
\bottomrule
\end{tabular}
\end{center}
The first file comes from P5's \src{lean/BigNumber/Definitions.lean}; the other
two come from the root BusyLean library. Their exact hashes are recorded in
\src{evidence/vendor_hashes.json}. No upstream compiled binaries or downloaded
font files are included.

\section{The unchanged object language}
\label{sec:language}

\subsection{Types, trees, and total interpretation}

The finite type codes satisfy
\[
 A,B ::= \mathsf{nat}\mid A\to B\mid\mathsf{wt}(A).
\]
Their denotations are natural numbers, function spaces, and the inductive
W-types
\[
 W(A)=\mathsf{leaf}\mid\mathsf{node}(A\to W(A)).
\]
Write $\W=W(\N)$. Raw terms have variables, typed abstractions, application,
zero, successor, natural recursion, and the leaf, node, and structural
recursion operations at all coded W-types. The existing interpreter is total
on raw syntax. It returns a type-tagged value and uses the original coercive
casts, including defaults for mismatches \cite{p5-defs}.

This last detail matters. We have not silently replaced the raw interpreter
with a different intrinsically typed language. In the numerical maximum,
non-numerical results score zero, exactly as in the existing approach. For the
particular comparison programs we prove that every application is correctly
typed, so the fallback behavior is never used in the successful computation.

A W-node's infinitely many children are represented by a function, not an
infinite list that must first be evaluated. Structural recursion on such a
node receives recursively defined functions on its children. The closed
construction only makes finitely many demands in any terminating natural
computation. The governing justification is the accepted structural recursion,
not an empirical observation that a few small runs finish.

\subsection{Levels, not de Bruijn indices}

The repository's variable lookup uses \emph{oldest-first levels}. An
abstraction prepends a new binding to the environment list, while level zero
continues to refer to the oldest binding. Consequently, adding local binders
does not change the level of a designated global interface variable.
The new seven-slot environment exploits this convention.

One must distinguish the list order from the level order: the environment
list is written newest first. Reversing that list incorrectly would make the
comparison program call entirely different values. The semantic proof and
reference interpreter explicitly test the intended convention.

\subsection{A symbolic rank bound for the enormous enumerations}

Let $\mathcal A_n$ and $\mathcal T_n$ denote the unchanged type and term lists
\code{allTypes n} and \code{allTerms n}. They are finite recursive lists. At
successor depth, the type list contains $\mathsf{nat}$, arrows between types
in the preceding list, and their W-types. The term list contains variables
with levels below the preceding depth, applications of preceding terms,
abstractions using preceding type annotations and bodies, and the primitive
recursors with preceding annotations.

Define type ranks by
\[
 \rank_A(\mathsf{nat})=1,\;
 \rank_A(A\to B)=1+\max(\rank_A(A),\rank_A(B)),\;
 \rank_A(\mathsf{wt}(A))=1+\rank_A(A).
\]
Term ranks are
\begin{align*}
 \rank(\mathsf{var}\,i)&=i+2,\\
 \rank(\mathsf{zero})=\rank(\mathsf{succ})&=1,\\
 \rank(\mathsf{app}(t,u))&=1+\max(\rank(t),\rank(u)),\\
 \rank(\mathsf{lam}(A,B,t))&=
   1+\max\{\rank_A(A),\rank_A(B),\rank(t)\},\\
 \rank(\mathsf{natrec}(A))&=1+\rank_A(A).
\end{align*}
The same one-plus-annotation rule applies to W-leaf and W-node constants;
W-recursion has rank one plus the maximum of its two annotation ranks.

\begin{lemma}[Enumeration inclusion]\label{lem:enumeration}
If $\rank_A(A)\le n$, then $A\in\mathcal A_n$. If $\rank(t)\le n$, then
$t\in\mathcal T_n$.
\end{lemma}
\begin{proof}
Use structural induction on the type or term, uniformly in the upper bound
$n$. The zero-depth cases are impossible because all ranks are positive. In
an application of rank at most $n+1$, both subterms have rank at most $n$;
the induction hypotheses put them in $\mathcal T_n$, and the product-list
membership lemma puts the application in $\mathcal T_{n+1}$. The abstraction
case additionally uses the type induction hypotheses. The variable inequality
$i+2\le n+1$ gives $i<n$, exactly the membership condition in the variable
prefix used for $\mathcal T_{n+1}$. Primitive terms follow from the singleton
and annotation-mapping parts of the enumerator. These arguments prove the
claim at every larger depth without expanding any particular large list.
\end{proof}

This is the route used in \src{Enumeration.lean}. Checking the proposed
comparison need not enumerate $\mathcal T_{42}$ or compute its length. The
small-depth reference counts in Section~\ref{sec:validation} are tests of the
implementation, not the method of proving the depth-42 membership fact.

\section{A fixed interface for changing interpretations}
\label{sec:interface}

\subsection{The model record}

A model is a triple of total operations
\begin{equation}\label{eq:model}
 M=(\rho_M,\sigma_M,\tau_M),
 \quad\begin{cases}
 \rho_M:\N\to\Fam\to\N\to\N,\\
 \sigma_M:\N\to\Fam\to\W\to\N\to\N,\\
 \tau_M:\N\to\W.
 \end{cases}
\end{equation}
The selector argument chooses a lower model. The fields represent its
reflection operation, forest recursion, and generated schedule, respectively.
A general model need not satisfy any coherence property for the elementary
finite-maximum lemmas. The constructors below supply coherent triples for
which the selected fields all come from the same lower interpretation.

The interface has a fixed finite type signature. This is important: simply
appending a new global constant at every iteration would make certificate
addresses and their ranks grow with the iteration count. Here the address
stays fixed and a numerical selector chooses the lower interpretation.
That selector is itself computed by a term of the object language.

\subsection{Seven stable global levels}

For a model $M$, external oracle $f:\Fam$, and input $n$, interpretation uses
the following oldest-first global levels:
\begin{center}
\begin{tabular}{@{}cll@{}}
\toprule
Level & Value & Object-language type\\
\midrule
0 & External oracle $f$ & $\N\to\N\to\N$\\
1 & Input $n$ & $\N$\\
2 & Inherited seed $\ps$ & $\N\to\N\to\N$\\
3 & Lower reflectors $\rho_M$ & $\N\to\Fam\to\N\to\N$\\
4 & Lower forests $\sigma_M$ & $\N\to\Fam\to\W\to\N\to\N$\\
5 & Lower schedules $\tau_M$ & $\N\to\W$\\
6 & Tower feedback $\Theta$ & $\N\to\N$\\
\bottomrule
\end{tabular}
\end{center}
All displayed types are abbreviations for finite type codes of the unchanged
language. The model record itself is a host-level organizing device; the
object language sees only these individually typed fields.

Write $\Eval_M(f,n,t)$ for the numerical score of the old raw interpreter
in this environment. The strengthened reflection operator is
\begin{equation}\label{eq:R}
 R_M(f,n)=n+1+
 \max\left\{\max_{i\le n} f(i,n),
              \max_{t\in\mathcal T_n}\Eval_M(f,n,t)\right\}.
\end{equation}
The forest operator is defined by structural recursion:
\begin{equation}\label{eq:S}
 \begin{split}
 S_M(g,\mathsf{leaf},n)&=R_M(g,n),\\
 S_M(g,\mathsf{node}(a),n)&=
       R_M\bigl(i\mapsto (m\mapsto S_M(g,a(i),m)),n\bigr).
 \end{split}
\end{equation}
The child is callable at an arbitrary dynamically computed numerical input;
we have not replaced the child family by a pointwise maximum.

\subsection{Generating schedules at all independent inputs}

Let $0_{\Fam}$ be the zero external oracle. Define
\[
 \td(M,m,t)=\operatorname{cast}_{\W}
     \bigl(\operatorname{interpret}(\operatorname{env}(M,0_{\Fam},m),t)\bigr).
\]
The cast includes the original default behavior when the resulting type is
not $\W$. Now set
\begin{equation}\label{eq:index}
 \idx(M)=\mathsf{node}\left(d\mapsto
  \mathsf{node}\left(m\mapsto
   \mathsf{node}\left(i\mapsto
       \td(M,m,\operatorname{nthD}(\mathcal T_d,i,\mathsf{zero}))
             \right)\right)\right).
\end{equation}
The three independent indices are syntax depth $d$, numerical parameter $m$,
and position $i$ in the finite term list. In particular, the construction
includes a program evaluated at any fixed input, not merely at an input
equal to its enumeration depth or position.

This is a definition by total function application and W-node construction.
It does not search over arbitrary programs for a proof of termination or
well-foundedness. Every denoted tree is already in the inductive type.

The number associated with a model and an input is
\begin{equation}\label{eq:value}
                 \Val(M,n)=S_M(\ps,\idx(M),n).
\end{equation}
The final contender will be one instance of this uniform construction.

\section{The depth-11 uniform certificate}
\label{sec:certificate}

\subsection{The certificate program}

For a selector term $s$, write the left-associated application
\begin{equation}\label{eq:call}
 \operatorname{call}(s)=
       \bigl(((v_4\ s)\ v_2)\ (v_5\ s)\bigr)\ v_1,
 \qquad
 A(s)=v_6\ \operatorname{call}(s),
\end{equation}
where $v_j$ denotes the variable at level $j$. Informally, the program selects
a lower forest, gives it the inherited seed and matching lower schedule,
evaluates at the current input, and applies tower feedback to the result.

\begin{lemma}[Exact rank]\label{lem:rank11}
For every selector $s$,
\[
 \rank(A(s))=\max\{11,\rank(s)+5\}.
\]
In particular, $\rank(A(\mathsf{zero}))=
\rank(A(v_1))=11$.
\end{lemma}
\begin{proof}
Let $r=\rank(s)$. The variable ranks of $v_4,v_5,v_6,v_2,v_1$ are
$6,7,8,4,3$. The first application $v_4\ s$ has rank
$\max(7,r+1)$; applying to $v_2$ gives $\max(8,r+2)$.
The other branch $v_5\ s$ has rank $\max(8,r+1)$, so the next application
has rank $\max(9,r+3)$. Applying to $v_1$ gives
$\max(10,r+4)$, and the final application of $v_6$ gives
$\max(11,r+5)$. Finally, zero has rank 1 and $v_1$ has rank 3.
\end{proof}

Thus a single small syntax tree certifies a comparison for all models and all
$n\ge11$. Its rank does not grow with the complexity of the model stored in
the environment.

\subsection{Why casts must be justified}

A correctly tagged application still uses the repository's \code{castValue}
function. At W-types, a self-cast traverses the tree with an identity branch
map, so it should not be dismissed as a purely definitional identity.

\begin{lemma}[Self-casts]\label{lem:casts}
For every finite type code $A$ and every $x\in\llbracket A\rrbracket$,
$\operatorname{castValue}(A,x)=x$.
\end{lemma}
\begin{proof}
First prove that the Boolean type equality test recognizes $A=A$, by type
induction. For every $w\in W(B)$, prove
$\operatorname{wmap}(\operatorname{id},w)=w$ by W-induction. The leaf case is
immediate. At a node, the induction hypotheses identify the child functions
pointwise; function extensionality identifies the functions themselves, and
congruence identifies the nodes.

Now induct on $A$ to show that the pair of casts from $A$ to itself is
$(\operatorname{id},\operatorname{id})$. The natural case is immediate. In
an arrow case, the recursively supplied domain and codomain casts are
identities, so both composite functions are identities. In a W-case, apply
the identity branch-map result. The desired statement is the first component
of this pair identity applied to $x$.
\end{proof}

The proposed Lean proof uses ordinary function extensionality, which is within
the repository's allowed axiom framework. Its actual transitive axiom set is
still a pending compiler audit. No equality of different tree notations or
ordinal representations is needed.

\begin{lemma}[Certificate semantics]\label{lem:certificate}
Suppose $s$ evaluates to the natural $k$ in the model environment. Then
\begin{equation}\label{eq:cert-eval}
 \Eval_M(f,n,A(s))=
     \Theta\bigl(\sigma_M(k,\ps,\tau_M(k),n)\bigr).
\end{equation}
For $s=\mathsf{zero}$ or $s=v_1$, this value is independent of the external
oracle $f$, with $k=0$ or $k=n$, respectively.
\end{lemma}
\begin{proof}
The level lookup gives exactly the values listed in Section~\ref{sec:interface}.
Every application in \eqref{eq:call} matches its declared domain and codomain.
Lemma~\ref{lem:casts} therefore removes the inserted self-casts. Applying the
selected forest to the seed, the selected schedule, and the input yields
$\sigma_M(k,\ps,\tau_M(k),n)$. Level 6 then applies $\Theta$. Neither zero
nor the input variable consults level 0, and no other part of this certificate
consults that level.
\end{proof}

\subsection{Lifting a witness through every forest}

\begin{lemma}[Strict reflection bounds]\label{lem:bounds}
For every model $M$, family $f$, and input $n$:
\begin{align*}
 n&<R_M(f,n),\\
 i\le n&\Longrightarrow f(i,n)<R_M(f,n),\\
 t\in\mathcal T_n&\Longrightarrow
             \Eval_M(f,n,t)<R_M(f,n).
\end{align*}
\end{lemma}
\begin{proof}
Every indicated value, except the already separately handled $n$, is at most
the maximum in \eqref{eq:R}. The positive summand $n+1$ makes the inequality
strict. Also $R_M(f,n)\ge n+1$ because its maximum is nonnegative.
The prefix and list maximum bounds are elementary inductions on the natural
prefix length and list membership.
\end{proof}

\begin{theorem}[Oracle-uniform witness lifting]\label{thm:lifting}
Fix $M,n,t$ with $t\in\mathcal T_n$. If a natural $b$ satisfies
\[
                 \forall f:\Fam,\quad\Eval_M(f,n,t)=b,
\]
then for every seed $g$ and every tree $a\in\W$,
\[
                         b<S_M(g,a,n).
\]
\end{theorem}
\begin{proof}
Split into the leaf and node cases of $a$. At a leaf, use
Lemma~\ref{lem:bounds} with $f=g$. At a node, use the same lemma with
$f(i,m)=S_M(g,a(i),m)$. The hypothesis explicitly quantifies over this family,
so the value is still $b$. No induction hypothesis or comparison between the
different oracles is required.
\end{proof}

Combining the preceding results yields, for every $M$ and $n\ge11$,
\begin{align}
 \Theta\bigl(\sigma_M(0,\ps,\tau_M(0),n)\bigr)&<\Val(M,n),
       \label{eq:zero-margin}\\
 \Theta\bigl(\sigma_M(n,\ps,\tau_M(n),n)\bigr)&<\Val(M,n).
       \label{eq:input-margin}
\end{align}
These two inequalities drive the entire numerical proof.

\begin{remark}[No oracle-monotonicity shortcut]
An object program can test whether an oracle result is zero; thus increasing
an oracle pointwise can decrease that program's output. This observation
alone does \emph{not} prove that the maximum defining $R_M$ is nonmonotone,
since a maximum can behave differently from an individual member. It does
show why pointwise comparison of oracle families cannot simply be pushed
through an arbitrary program. Our proofs make no such inference about either
programs or reflectors. They use a fixed oracle-independent witness instead.
\end{remark}

\section{Autonomous model ramification}
\label{sec:models}

\subsection{One step and all finite steps}

Define the successor model by exposing the current model's operations:
\begin{equation}\label{eq:next}
 \begin{split}
 \rho_{\Next(M)}(k)&=R_M,\\
 \sigma_{\Next(M)}(k)&=S_M,\\
 \tau_{\Next(M)}(k)&=\idx(M).
 \end{split}
\end{equation}
The selector is ignored at this one-step interface. Let
\[
 M^{(0)}=M,\qquad M^{(k+1)}=\Next(M^{(k)}).
\]
The closure under all finite steps, denoted $\Omega(M)$, has fields
\begin{equation}\label{eq:omega}
 \rho_{\Omega(M)}(k)=R_{M^{(k)}},\quad
 \sigma_{\Omega(M)}(k)=S_{M^{(k)}},\quad
 \tau_{\Omega(M)}(k)=\idx(M^{(k)}).
\end{equation}
The notation $\Omega$ is a convenient name for this operator. It is
\emph{not} an asserted ordinal notation or an identification with a set-theoretic
ordinal. The defining recursion for $M^{(k)}$ is ordinary natural recursion.

A program can calculate a selector $k$ much larger than its own input or
syntax depth. The specific lower bound will only need selectors zero and
$n$, but the full operator retains more computational capability.

\subsection{Retaining a base and every child}

For a base model $B$ and a family $i\mapsto M_i$, let
$J=\MJoin(B,(M_i)_i)$ expose the operations of $B$ at selector zero and of
$M_i$ at selector $i+1$:
\begin{equation}\label{eq:join}
 \begin{array}{lll}
 \rho_J(0)=R_B,&\sigma_J(0)=S_B,&\tau_J(0)=\idx(B),\\
 \rho_J(i+1)=R_{M_i},&\sigma_J(i+1)=S_{M_i},&
                         \tau_J(i+1)=\idx(M_i).
 \end{array}
\end{equation}
The base is retained even when there are children. This explicit zero slot is
what makes the uniform base lower bound independent of which schedule is
used. Each child is retained as a full interpreter interface, rather than as
one numerical output.

Define ramification by structural recursion on $a\in\W$:
\begin{equation}\label{eq:ramify}
 \begin{split}
 \Ram(B,\mathsf{leaf})&=\Omega(B),\\
 \Ram(B,\mathsf{node}(a))&=
    \Omega\!\left(\MJoin\bigl(B,(\Ram(B,a(i)))_i\bigr)\right).
 \end{split}
\end{equation}
Thus every node collects all the recursively constructed child models and
closes the resulting interface under all finite further steps. This recursion
has values in the fixed host type \code{Model}. It does not quantify over all
host-level definitions and does not require an internal truth predicate.

\subsection{The closed final model}

The initial model $M_0$ exposes the existing components:
\begin{equation}\label{eq:base}
 \rho_{M_0}(k)=R_{\mathrm{old}},\qquad
 \sigma_{M_0}(k)=S_{\mathrm{old}},\qquad
 \tau_{M_0}(k)=I_{\mathrm{old}}.
\end{equation}
The closed schedule and final model are
\begin{equation}\label{eq:final}
 \begin{aligned}
 M_\omega&=\Omega(M_0),\\
 a_*&=\idx(M_\omega),\\
 M_*&=\Ram(M_\omega,a_*),\\
 M_{\mathrm{final}}&=\Omega(M_*),\\
 C_7&=\Val(M_{\mathrm{final}},42).
 \end{aligned}
\end{equation}
The defining dependency graph has the forward order
\[
 M_0\ \longrightarrow\ M_\omega\ \longrightarrow\ a_*
 \ \longrightarrow\ M_*\ \longrightarrow\ M_{\mathrm{final}}
 \ \longrightarrow\ C_7.
\]
The schedule is generated by an \emph{already defined lower} model. The
ramified model does not occur in the definition of its own schedule. The term
``autonomous'' refers to the language-generated choice of schedule, not to a
circular equation for a universal self-interpreter.

\section{The complete quantitative proof}
\label{sec:quantitative}

\subsection{Elementary facts about tower feedback}

\begin{lemma}\label{lem:theta-properties}
The function $\Theta$ is monotone on $\N$, and $x<\Theta(x)$ for every $x$.
Every iterate $\Iter\Theta k$ is monotone and inflationary.
\end{lemma}
\begin{proof}
First, $2^x>x$ for every natural $x$. At zero this is $1>0$; the successor
step follows from $2^{x+1}=2\cdot2^x$ and $2^x\ge x+1$.
The function $E(x)=2^x$ is monotone. Its iterates are therefore monotone in
the starting value by induction, and inflationary because each individual
step is inflationary. They are also monotone in the iteration count.

If $x\le y$, apply these two monotonicity facts to the initial values
$x+1\le y+1$ and counts $x+1\le y+1$ to obtain
$\Theta(x)\le\Theta(y)$. Furthermore, the initial value $x+1$ is preserved or
increased by its $x+1$ iterations, so $x<x+1\le\Theta(x)$.
Induction gives monotonicity and inflationarity of every $\Theta$ iterate.
\end{proof}

For clarity, none of the comparison arguments below requires strict
monotonicity of $\Theta$. Ordinary monotonicity plus the strict inequalities
from reflection suffices.

\subsection{Successor, finite, and omega bounds}

\begin{proposition}[Model amplification]\label{prop:amplification}
For every model $M$ and $n\ge11$,
\begin{align}
 \Theta(\Val(M,n))&<\Val(\Next(M),n),\label{eq:next-bound}\\
 \Iter\Theta k(\Val(M,n))&\le\Val(M^{(k)},n)
          &&(k\in\N),\label{eq:finite-bound}\\
 \Theta(\Val(M^{(n)},n))&<\Val(\Omega(M),n),\label{eq:omega-diag}\\
 \Iter\Theta{n+1}(\Val(M,n))&<\Val(\Omega(M),n),\label{eq:omega-bound}\\
 \Theta(\Val(M,n))&<\Val(\Omega(M),n).\label{eq:omega-base}
\end{align}
\end{proposition}
\begin{proof}
Use the zero-selector certificate \eqref{eq:zero-margin} in the model
$\Next(M)$. By its field definitions, the selected forest and schedule are
$S_M$ and $\idx(M)$, so this is exactly \eqref{eq:next-bound}.

For \eqref{eq:finite-bound}, induct on $k$. The base case is equality. Given
the bound at $k$, monotonicity of $\Theta$ gives
\[
 \Iter\Theta{k+1}(\Val(M,n))
 \le\Theta(\Val(M^{(k)},n))
 <\Val(M^{(k+1)},n).
\]
Dropping strictness proves the desired successor inequality.

Use the input-selector certificate \eqref{eq:input-margin} in $\Omega(M)$.
Its selector $n$ exposes $M^{(n)}$, giving \eqref{eq:omega-diag}. Applying
$\Theta$ to \eqref{eq:finite-bound} with $k=n$ and composing proves
\eqref{eq:omega-bound}. Finally, using the zero-selector certificate instead
exposes $M^{(0)}=M$ and proves \eqref{eq:omega-base} directly.
\end{proof}

\begin{proposition}[Preservation through joins and ramification]
\label{prop:ram-base}
For every base $B$, every child family $(M_i)$, every tree $a$, and $n\ge11$,
\begin{align}
 \Theta(\Val(B,n))&<\Val(\MJoin(B,(M_i)),n),\label{eq:join-base}\\
 \Theta(\Val(B,n))&<\Val(\Ram(B,a),n).\label{eq:ram-base}
\end{align}
\end{proposition}
\begin{proof}
The zero slot of the join contains precisely $S_B$ and $\idx(B)$. Thus
\eqref{eq:join-base} is another instance of the zero-selector certificate.
For \eqref{eq:ram-base}, split into the tree's two constructors. At a leaf,
use \eqref{eq:omega-base}. At a node, write
$J=\MJoin(B,(\Ram(B,a(i)))_i)$. Then
\[
 \Theta(\Val(B,n))<\Val(J,n)
       <\Theta(\Val(J,n))<\Val(\Omega(J),n).
\]
The first inequality is the join bound, the middle one is
Lemma~\ref{lem:theta-properties}, and the last is the omega base bound. The
rightmost value is the required ramification value. No inductive comparison
of the children is needed for this particular lower bound.
\end{proof}

\subsection{From the old family to the final model}

By the fields of $M_0$, the zero-selector certificate reconstructs the old
family exactly. For every $n\ge11$,
\begin{equation}\label{eq:base-bound}
                     \Theta(\FF_6(n))<\Val(M_0,n).
\end{equation}
This is an equality-of-computation comparison, not an ordinal estimate.

\begin{theorem}[Uniform separation]\label{thm:uniform}
For every natural $n\ge11$,
\begin{align}
 \Iter\Theta{n+2}(\FF_6(n))&<\Val(M_\omega,n),\label{eq:initial-final}\\
 \Iter\Theta{n+3}(\FF_6(n))&<\Val(M_*,n),\label{eq:ram-final}\\
 \Iter\Theta{n+4}(\FF_6(n))&<\Val(M_{\mathrm{final}},n).
 \label{eq:final-final}
\end{align}
\end{theorem}
\begin{proof}
Apply the monotone function $\Iter\Theta{n+1}$ to the non-strict consequence of
\eqref{eq:base-bound}, then use \eqref{eq:omega-bound}:
\[
 \Iter\Theta{n+2}(\FF_6(n))
 \le\Iter\Theta{n+1}(\Val(M_0,n))
 <\Val(\Omega(M_0),n).
\]
This proves \eqref{eq:initial-final}. Apply $\Theta$ to its non-strict
consequence and then use \eqref{eq:ram-base} with base $M_\omega$ and tree
$a_*$. The result is \eqref{eq:ram-final}. Apply $\Theta$ once more and use
\eqref{eq:omega-base} with $M=M_*$ to get \eqref{eq:final-final}.
\end{proof}

Theorem~\ref{thm:main} now follows by taking $n=42$, using
$C_6=\FF_6(42)$ and the definition of $C_7$. Inflationarity gives
$C_6\le\Iter\Theta{46}(C_6)$, so the ordinary strict comparison follows too.

The proof has deliberately not used most of the strength of the final omega
closure or the structure of $a_*$. It establishes a readily inspectable
numerical margin first. The next section explains a structural guarantee that
would be lost if the schedule were replaced by an arbitrary single leaf.

\section{Cofinality of the generated ramification schedule}
\label{sec:cofinality}

\subsection{Selecting a particular child}

Let $\overline k$ be the unary numeral obtained from zero by $k$ applications
of successor. Then $\rank(\overline k)=k+1$, and
\begin{equation}\label{eq:numeral-rank}
                  \rank(A(\overline k))=\max\{11,k+6\}.
\end{equation}
The use of unary numerals is a proof convenience, not a claim that it gives
the most compact selectors. The full language can compute much larger
selectors using shorter programs.

\begin{lemma}[Visible child amplification]\label{lem:child}
Let $a:\N\to\W$. If $n\ge11$ and $i+7\le n$, then
\[
 \Theta\bigl(\Val(\Ram(B,a(i)),n)\bigr)
       <\Val(\Ram(B,\mathsf{node}(a)),n).
\]
In particular the same inequality holds with $\Theta$ omitted from the left.
\end{lemma}
\begin{proof}
Set $M_i=\Ram(B,a(i))$ and $J=\MJoin(B,(M_i))$. The term
$A(\overline{i+1})$ has rank $\max(11,i+7)\le n$, independently of the
external oracle. The selected fields of $J$ are $S_{M_i}$ and $\idx(M_i)$.
The uniform witness lemma gives
\[
                   \Theta(\Val(M_i,n))<\Val(J,n).
\]
The omega base bound gives
$\Val(J,n)<\Theta(\Val(J,n))<\Val(\Omega(J),n)$.
The last model is exactly $\Ram(B,\mathsf{node}(a))$.
Inflationarity of $\Theta$ gives the unamplified consequence.
\end{proof}

\subsection{Every fixed definable schedule is eventually dominated}

\begin{theorem}[Interpretation-relative cofinality]\label{thm:cofinality}
For every base model $B$, model $M$, raw term $t$, and natural parameter $m$,
there is a natural $N$ such that for every $n\ge N$,
\begin{equation}\label{eq:cofinality}
 \Val\bigl(\Ram(B,\td(M,m,t)),n\bigr)
       <\Val\bigl(\Ram(B,\idx(M)),n\bigr).
\end{equation}
\end{theorem}
\begin{proof}
Let $d=\rank(t)$. Lemma~\ref{lem:enumeration} gives
$t\in\mathcal T_d$. A simple induction on list membership supplies some
position $i$ for which
$\operatorname{nthD}(\mathcal T_d,i,\mathsf{zero})=t$.
Choose
\begin{equation}\label{eq:threshold}
                N=\max\{11,d+7,m+7,i+7\}.
\end{equation}
For $n\ge N$, apply the unamplified conclusion of
Lemma~\ref{lem:child} three times: first to position $i$, then to parameter
$m$, and finally to depth $d$. The three node layers are exactly those in
\eqref{eq:index}. Substituting the list-position equality at the innermost
leaf gives \eqref{eq:cofinality}.
\end{proof}

For the final construction, take $B=M=M_\omega$. The intermediate model
$M_*$ eventually dominates the model ramification along every fixed tree
that the initially omega-strengthened interpretation produces, at every
fixed numerical parameter. The final omega closure further dominates this
intermediate model at all sufficiently large inputs.

Several scope restrictions are essential. The threshold can be large because
the position in a finite enumeration can be large. The theorem therefore does
not say that every arbitrarily complicated schedule is already dominated at
input 42. It is relative to a specified model's denotation of raw terms, not
to the class of all Lean-definable trees or all computable well-orders. It is
a numerical comparison of the resulting ramification constructions, not an
assertion that the index has been assigned a particular ordinal rank.

\section{Diagonal escape without a halting oracle}
\label{sec:diagonal}

\begin{theorem}[No eventual self-majorant]\label{thm:diagonal}
Fix a model $M$, a raw term $t$, and any input-dependent family of external
oracles $n\mapsto f_n$. There is no $N$ such that
\[
 \forall n\ge N,\qquad R_M(f_n,n)\le\Eval_M(f_n,n,t).
\]
\end{theorem}
\begin{proof}
Suppose such an $N$ exists and let $n=\max\{N,\rank(t)\}$. By
Lemma~\ref{lem:enumeration}, $t\in\mathcal T_n$. The strict program bound of
Lemma~\ref{lem:bounds} gives
\[
 \Eval_M(f_n,n,t)<R_M(f_n,n)\le\Eval_M(f_n,n,t),
\]
a contradiction.
\end{proof}

This diagonal theorem concerns a \emph{fixed} interpretation with its fixed
interface. The successor construction can denote computations associated with
the lower interpretation because it changes that interface. Consequently,
there is no contradiction between the diagonal theorem and the successful
short certificate that calls the lower forest in the next model.

Nor is there an appeal to a Busy Beaver maximum. The finite lists contain
terms of an independently total interpreted language. We neither filter all
Turing programs by whether they halt nor assign values using a noncomputable
halting predicate. The diagonal quantifiers range over the syntax of the
specified language, not over every possible terminating host-language program.

\section{Computability and foundational discipline}
\label{sec:computability}

\subsection{Why the closed number is computable in principle}

The proposed computation uses the following forms of recursion. The existing
interpreter recurses structurally on raw syntax, with natural and W-recursion
provided at the denoted types. Prefix maxima and finite term lists use natural
or list recursion. The new forest operator recurses structurally on its
W-tree. The finite model tower uses natural recursion, and model ramification
uses W-tree recursion. All later expressions in \eqref{eq:final} are finite
compositions of these operations.

This is a constructive recipe, not a classical existence definition of a
natural. Once the definitions are accepted by the pinned Lean checker, their
termination is supported by the same inductive-recursion discipline used by
the baseline \cite{lean-recursion,lean-inductive}. Independently of compiler
acceptance, the displayed mathematical recursions give the total construction
and its proofs; the exact Lean source's acceptance is the separate validation
question recorded in this package.

For a generic argument $M$, the mathematical model record permits arbitrary
total functions, which need not be computable. No claim is made that
$\Val(M,n)$ is uniformly computable given a noncomputable external model.
The computability assertion is about the \emph{closed} $M_0$ and the explicit
constructors applied to it. Every model field in that closed dependency graph
has an algorithmic definition; no noncomputable data are supplied at a leaf.

There is no practical route to writing down the final decimal expansion.
Computability in principle is not a resource-feasibility claim. Even the
stated lower bound is beyond physical evaluation, and executing the finite
maximum at depth 42 would itself be inappropriate. The build and proof
strategy is to check symbolic definitions and inequalities while never
normalizing either contender.

\subsection{No hidden universes or paradoxical self-evaluation}

The new record has a fixed collection of ordinary higher-order function
fields. Iterating its construction does not require extending the object type
syntax with a universe of all models, nor with a type of all host programs.
The same finite type codes interpret each field at every stage. The host
language defines the family of interpretations using its own structural
recursion.

Similarly, $\idx(M)$ is not defined as a join of all trees expressible using
$\idx(M)$ as a new primitive. It is defined using the operations \emph{already
stored in} $M$. Only after that index has been defined can \code{next M}
expose it. This sequencing is exactly what keeps the diagonal argument and
the construction compatible.

The use of extensional equality in the cast proof is explicit. There is no
\code{Classical.choice}, admitted theorem, custom axiom, or unsafe evaluation
mechanism intentionally present in the new source. This source-level design
statement is weaker than a completed transitive axiom audit. The latter must
be obtained from the compiled environment, not from a grep result.

\subsection{What has and has not been measured by ordinals}

The upstream review describes rough ordinal calibrations for the language and
its proposals, and labels them informal \cite{repo-review}. We do not use
those calibrations as premises. A proof-theoretic ordinal of a theory, an
ordinal notation system representing tree schedules, and a finite numerical
comparison at input 42 are different mathematical objects. Identifying one
with another requires a precisely stated translation and quantitative control
of fundamental sequences or coding costs.

The present contribution avoids this gap. It gives an exact program witness,
an exact rank cutoff, exact iteration inequalities, and an explicit eventual
cofinality threshold. It is meaningful even without determining which
ordinal-analysis framework best describes the overall model hierarchy.

\section{Formalization architecture and actual validation}
\label{sec:validation}

\subsection{Source modules}

The new library consists of 823 lines across eight modules, with 32
\code{def} declarations and 84 \code{theorem} declarations according to the
included lexical audit. The declaration counts exclude structures and
abbreviations and are not a proof-complexity measure.

\begin{center}
\begin{tabular}{@{}p{0.28\textwidth}p{0.66\textwidth}@{}}
\toprule
Module & Mathematical role\\
\midrule
\src{Definitions.lean} & Fixed model interface, new reflection and forest,
  generated index, model constructors, closed contender.\\
\src{Arithmetic.lean} & Inflationarity and monotonicity of exponential
  feedback and its iterations.\\
\src{Enumeration.lean} & Type/term rank inclusion and the depth-11
  certificates, without expanding large enumerations.\\
\src{Semantics.lean} & Self-cast identity, exact typed application,
  certificate evaluation.\\
\src{Reflection.lean} & Finite maximum inequalities, oracle-uniform
  witness lifting, diagonal non-majorization.\\
\src{Comparison.lean} & Successor/omega/ramification bounds and the
  46-fold comparison with the pinned contender.\\
\src{Selectors.lean} & Unary selectors and the explicit child
  visibility cost.\\
\src{Cofinality.lean} & Every fixed definable schedule is eventually
  dominated by the generated ramification schedule.\\
\bottomrule
\end{tabular}
\end{center}

The actual proposed scalar theorem names are
\begin{lstlisting}
BusyLean.contender_6_tower_lt_contender_7
BusyLean.contender_6_lt_contender_7
\end{lstlisting}
The strongest uniform numerical theorem is \code{BusyLean7.final\_bound};
the structural theorem is
\code{BusyLean7.ramify\_index\_eventually\_dominates}.
The package's root \src{BusyLean7.lean} imports both comparison and cofinality.

\subsection{Independent finite tests}

The Python reference implementation mirrors the raw syntax constructors,
finite type denotations, casting convention, oldest-first lookup, and small
finite enumerations. It has explicit safeguards against accidentally
attempting large computations. The test run recorded in
\src{evidence/python_tests.json} passed 1,830 assertions.

\begin{center}
\begin{tabular}{@{}lrrrrr@{}}
\toprule
Depth & 0 & 1 & 2 & 3 & 4\\
\midrule
Number of raw terms & 0 & 2 & 13 & 308 & 147,129\\
Mock-model numerical maximum & 0 & 0 & 1 & 3 & 5\\
\bottomrule
\end{tabular}
\end{center}

The second row of numbers concerns a deliberately small \emph{mock model};
it is not a sequence of actual values of $C_6$, $C_7$, or the proposed
self-ramified model. The tests also include 976 checks of the comparison
certificate across different numerical inputs, external oracles, and
selectors, as well as 48 small reflector checks. The tested tower values are
$2$, $16$, and $2^{256}$. Evaluation beyond the guarded tiny cases is refused.

These checks are useful for finding level-order errors, wrong application
nesting, incorrect syntax-rank offsets, and mistakes in the fallback cast
behavior. They cannot establish the universal claims of the article and
cannot establish that Lean accepts a proof term. In particular, a passing
independent interpreter cannot serve as a substitute for the Lean kernel.

\subsection{The audit ledger}

\begin{center}
\begin{tabular}{@{}p{0.47\textwidth}p{0.47\textwidth}@{}}
\toprule
Check & Recorded outcome\\
\midrule
Pinned baseline Git blob identities & All three matched exactly.\\
Finite Python reference tests & Passed; 1,830 assertions.\\
Lexical scan of new computational sources & No prohibited tokens found;
  no direct reference to the forbidden named winners.\\
New Lean definition elaboration & Not run; compiler unavailable.\\
New Lean proof checking & Not run; compiler unavailable.\\
Transitive axiom/dependency audit & Supplied, but not run.\\
15-second / 60-second clean-build budgets & Unverified.\\
Evaluation of either contender & Not attempted.\\
\bottomrule
\end{tabular}
\end{center}

A lexical scan is not a parser, a termination checker, or a dependency-closure
computation. The proposed Lean audit separately traverses constants in both
types and values, prints the axioms of submitted roots, and tests forbidden
constant reachability. Its design follows the pinned repository's audit
pattern \cite{repo-audit}; its execution remains outstanding.

No \code{sorry} or custom axiom is inserted to make an apparent build succeed.
If the first real compiler run reports an error, that should be repaired and
recorded honestly rather than suppressed. The source is modular precisely so
that enumeration, semantic casts, numerical comparisons, and cofinality can
be checked in that order without evaluating a huge natural.

\section{Reproduction and integration}
\label{sec:reproduction}

The ZIP is standalone for the proposed new definition and direct comparison:
the three necessary old definition files are vendored. It does not require
fetching the full repository or Mathlib merely to inspect or attempt the new
build. The toolchain itself must be installed separately.

From the package root, the available finite checks are
\begin{lstlisting}
python3 python/test_reference.py
python3 scripts/source_audit.py
\end{lstlisting}
With Lean 4.34.0 and Lake installed, run
\begin{lstlisting}
python3 scripts/check.py
\end{lstlisting}
The script checks the toolchain, removes the package's Lean build directory,
builds the definition phase and proof phase separately, runs safe smoke tests,
and runs the axiom/dependency audit. Missing Lean yields a non-success result
and an explicit \code{not\_run} receipt, not a fabricated pass. The script
also records whether measured definition and proof times meet the repository
budgets. A plain attempted library build is
\begin{lstlisting}
cd lean
lake build BusyLean7
\end{lstlisting}

The optional \src{integration/Transitive.lean} imports the full pinned
repository's existing \src{BusyLean/Comparison.lean} and composes its two
strict inequalities with the proposed $C_6<C_7$ theorem. That optional file is
outside the standalone build because the old comparison has additional
port-agreement dependencies that are not needed for the new direct bound.
It provides the intended statements
\begin{lstlisting}
BusyLean.contender_5_lt_contender_7
BusyLean.p5_lt_contender_7
\end{lstlisting}
without modifying or redefining the old contender constants.

For integration into the original repository, copy only the new namespace
modules and the desired tests, retaining the pinned baseline references for
review. The old files in this ZIP are source witnesses, not proposed changes
to their upstream contents. See \src{docs/INTEGRATION.md} for the concrete
layout and validation order.

Build this article with
\begin{lstlisting}
bash scripts/build_pdf.sh
\end{lstlisting}
The script runs three \code{pdflatex} passes and writes build logs to the
evidence directory. The delivered PDF was compiled from the accompanying
source; its validation receipt is separate from the Lean receipt.

\section{Further research questions}
\label{sec:future}

\subsection{Complete kernel certification before extending the hierarchy}

The immediate priority is a Lean 4.34.0 build of all eight modules, followed
by the transitive axiom/dependency audit and the clean timing measurement.
The generic mathematical lemmas have intentionally been separated from
scalar specialization to reduce accidental normalization. Any required
elaboration repairs should preserve the defining equations and the exact
rank-11 certificate, with a new validation receipt explaining the changes.
Only after that pass should the proposal be described as a certified
contest entry.

\subsection{Certify the other proposal constants by translations}

An exact $C_7>P_j$ statement for each of P1--P4 requires a concrete comparison,
not reliance on the repository's heuristic ranking. One route is to encode
their primitive tree operations and iteration functionals as fixed programs
in a sufficiently low model, prove semantic agreement, and track syntax
rank and parameter overhead. Another is a direct uniform interpretation of
each proposal's definition into the current forest interface. The finite
input 42 makes overhead accounting essential; eventual domination by itself
does not settle the scalar inequality.

\subsection{Reduce selection overhead without changing semantics}

The threshold $i+7\le n$ in the cofinality proof comes from a unary numeral.
Replacing it with a certified binary or arithmetic construction should make
many branches visible at much smaller input. The correct theorem would
state the cost in terms of the length of a verified selector program, not in
terms of the child index itself. A general oracle-uniform selector lemma is
already present conceptually; extending its library is a natural next step.
The three independent schedule indices should be retained during such a
compression.

\subsection{Use more of the final model in the numerical lower bound}

The final proof only uses the zero slot of the outermost omega closure and a
base-preservation bound through the generated schedule. Stronger explicit
bounds should exploit the input selector there, and then use identified
branches of $a_*$ with small certified positions. The challenge is not
merely obtaining a larger expression, but doing so with a transparent
comparison theorem that avoids implicit monotonicity assumptions about
interpreters.

\subsection{Develop a precise ordinal analysis separately}

A systematic calibration could define a notation system for the model
constructors, prove a simulation theorem relating that system to the W-tree
semantics, and derive upper and lower bounds for numerical evaluations.
Each bridge should specify its coding costs and parameter shifts. It would
be inappropriate to identify $\Omega(M)$ with a familiar ordinal operation
solely because both involve a countable family. The exact results in this
article can serve as constraints and test cases for such an analysis.

\subsection{Generalize the interface and retain quantitative witnesses}

The fixed-signature method applies whenever a total language can expose
families of lower evaluators and their recursion schedules. Higher-type
selectors, richer strictly positive data types, or multiple interacting
schedule families could produce stronger variants. The key design target is
an invariant short witness whose semantic value is independent of the
current external oracle. Without such a witness, plausible structural
strengthening can become difficult to compare at a fixed natural input.

\section{Conclusion}

The new number is obtained from a hierarchy of interpreters rather than from
an unexplained enormous numeral or a halting oracle. A fixed interface makes
lower reflection, forest recursion, and language-generated schedules
available to the next interpretation. Natural iteration and structural
model-valued W-recursion organize those interpretations without circular
self-reference.

A single depth-11 certificate proves a large, explicit, uniform gap. The
result at the repository input is $\Iter\Theta{46}(C_6)<C_7$. The generated
schedule additionally satisfies an exact eventual domination theorem for
ramification along every fixed tree definable in the lower model. These are
proved mathematical statements about the displayed construction, not
heuristic ordinal rankings.

The remaining boundary is computational verification of the submitted Lean
text. The package states that boundary prominently, provides reproducible
checking and audit scripts, and records only checks actually performed. This
keeps the substantial mathematical proposal separate from the certification
work that still has to be completed.

\appendix
\section{Complete new computational definition}
\label{app:code}

The following is the actual definition module included in the package. It
imports the pinned current definition to reuse the generic old components;
its new contender body does not reference the old contender or final unary
family. Proofs and the transitive dependency audit are in separate files.

\lstinputlisting[style=lean,basicstyle=\ttfamily\scriptsize]{../lean/BusyLean7/Definitions.lean}

\section{Correspondence between mathematical statements and sources}
\label{app:map}

\begin{longtable}{@{}p{0.40\textwidth}p{0.54\textwidth}@{}}
\toprule
Statement & Proposed Lean declaration\\
\midrule
\endhead
Enumeration inclusion & \code{BusyLean7.term\_mem}\\
Identity W-map & \code{BusyLean7.wmap\_identity}\\
Identity self-cast & \code{BusyLean7.castValue\_self}\\
Certificate semantics & \code{BusyLean7.interpret\_amplified}\\
Successor margin & \code{BusyLean7.next\_margin}\\
Finite iteration bound & \code{BusyLean7.finite\_bound}\\
Omega diagonal margin & \code{BusyLean7.omega\_diagonal\_margin}\\
Omega iteration bound & \code{BusyLean7.omega\_bound}\\
Ramification base margin & \code{BusyLean7.ramify\_margin}\\
Uniform final bound & \code{BusyLean7.final\_bound}\\
46-fold scalar bound & \code{BusyLean.contender\_6\_tower\_lt\_contender\_7}\\
Strict scalar comparison & \code{BusyLean.contender\_6\_lt\_contender\_7}\\
Eventual schedule domination & \code{BusyLean7.ramify\_index\_eventually\_dominates}\\
Diagonal non-majorization & \code{BusyLean7.no\_eventual\_self\_majorant}\\
\bottomrule
\end{longtable}

\noindent All entries in this correspondence table refer to supplied source,
not to completed Lean kernel checks.

\begin{thebibliography}{99}
\bibitem{repo-readme}
BusyLean contributors, \emph{BusyLean: name the biggest number, in Lean},
root README at commit
\code{86b360648163f0249eab79539438b6a83dfe8215}, retrieved October 12, 2026.
\href{https://github.com/VladimirReshetnikov/BusyLean/blob/86b360648163f0249eab79539438b6a83dfe8215/README.md}{Pinned README and contest rules}.

\bibitem{repo-c6}
BusyLean contributors, \emph{Contender 6},
\src{BusyLean/Contender6.lean}, same pinned commit.
\href{https://github.com/VladimirReshetnikov/BusyLean/blob/86b360648163f0249eab79539438b6a83dfe8215/BusyLean/Contender6.lean}{Pinned definition}.

\bibitem{repo-reflection}
BusyLean contributors, \emph{Reflection over the W-language},
\src{BusyLean/Reflection.lean}, same pinned commit.
\href{https://github.com/VladimirReshetnikov/BusyLean/blob/86b360648163f0249eab79539438b6a83dfe8215/BusyLean/Reflection.lean}{Pinned reflection source}.

\bibitem{p5-defs}
BusyLean contributors, \emph{A conservative finite-type W-language and a
constructive large-number entry}, proposal P5,
\src{lean/BigNumber/Definitions.lean}, same pinned commit.
\href{https://github.com/VladimirReshetnikov/BusyLean/blob/86b360648163f0249eab79539438b6a83dfe8215/proposals/P5-w-reflection/lean/BigNumber/Definitions.lean}{Pinned language and interpreter}.

\bibitem{repo-comparison}
BusyLean contributors, \emph{Contender 6: comparison theorems},
\src{BusyLean/Comparison.lean}, same pinned commit.
\href{https://github.com/VladimirReshetnikov/BusyLean/blob/86b360648163f0249eab79539438b6a83dfe8215/BusyLean/Comparison.lean}{Pinned comparison theorems}.

\bibitem{repo-review}
BusyLean contributors, \emph{Review of the five incoming proposals},
\src{proposals/REVIEW.md}, same pinned commit.
\href{https://github.com/VladimirReshetnikov/BusyLean/blob/86b360648163f0249eab79539438b6a83dfe8215/proposals/REVIEW.md}{Pinned review; its ordinal calibrations are explicitly informal}.

\bibitem{repo-audit}
BusyLean contributors, \emph{Audit}, \src{test/Audit.lean}, same pinned commit.
\href{https://github.com/VladimirReshetnikov/BusyLean/blob/86b360648163f0249eab79539438b6a83dfe8215/test/Audit.lean}{Pinned dependency and axiom audit pattern}.

\bibitem{lean-recursion}
Lean development team, \emph{The Lean Language Reference: Recursive
Definitions}, official documentation, consulted October 12, 2026.
\href{https://lean-lang.org/doc/reference/latest/Definitions/Recursive-Definitions/}{Recursive definitions and termination checking}.

\bibitem{lean-inductive}
Lean development team, \emph{The Lean Language Reference: Inductive Types},
official documentation, consulted October 12, 2026.
\href{https://lean-lang.org/doc/reference/latest/The-Type-System/Inductive-Types/}{Inductive types and their recursors}.

\bibitem{lean-release}
Lean development team, \emph{Lean 4.34.0 release}, official release record.
\href{https://github.com/leanprover/lean4/releases/tag/v4.34.0}{Pinned toolchain release}.
\end{thebibliography}

\end{document}
